% Options for packages loaded elsewhere
\PassOptionsToPackage{unicode}{hyperref}
\PassOptionsToPackage{hyphens}{url}
\documentclass[
 12pt,
]{ctexart}
\usepackage{xcolor}
\usepackage{amsmath,amssymb}
\setcounter{secnumdepth}{-\maxdimen} % remove section numbering
\usepackage{iftex}
\ifPDFTeX
 \usepackage[T1]{fontenc}
 \usepackage[utf8]{inputenc}
 \usepackage{textcomp} % provide euro and other symbols
\else % if luatex or xetex
 \usepackage{unicode-math} % this also loads fontspec
 \defaultfontfeatures{Scale=MatchLowercase}
 \defaultfontfeatures[\rmfamily]{Ligatures=TeX,Scale=1}
\fi
\usepackage{lmodern}
\ifPDFTeX\else
 % xetex/luatex font selection
\fi
% Use upquote if available, for straight quotes in verbatim environments
\IfFileExists{upquote.sty}{\usepackage{upquote}}{}
\IfFileExists{microtype.sty}{% use microtype if available
 \usepackage[]{microtype}
 \UseMicrotypeSet[protrusion]{basicmath} % disable protrusion for tt fonts
}{}
\makeatletter
\@ifundefined{KOMAClassName}{% if non-KOMA class
 \IfFileExists{parskip.sty}{%
 \usepackage{parskip}
 }{% else
 \setlength{\parindent}{0pt}
 \setlength{\parskip}{6pt plus 2pt minus 1pt}}
}{% if KOMA class
 \KOMAoptions{parskip=half}}
\makeatother
\usepackage{longtable,booktabs,array}
\usepackage{calc} % for calculating minipage widths
% Correct order of tables after \paragraph or \subparagraph
\usepackage{etoolbox}
\makeatletter
\patchcmd\longtable{\par}{\if@noskipsec\mbox{}\fi\par}{}{}
\makeatother
% Allow footnotes in longtable head/foot
\IfFileExists{footnotehyper.sty}{\usepackage{footnotehyper}}{\usepackage{footnote}}
\makesavenoteenv{longtable}
\setlength{\emergencystretch}{3em} % prevent overfull lines
\providecommand{\tightlist}{%
 \setlength{\itemsep}{0pt}\setlength{\parskip}{0pt}}
\usepackage{bookmark}
\IfFileExists{xurl.sty}{\usepackage{xurl}}{} % add URL line breaks if available
\urlstyle{same}
\hypersetup{
 hidelinks,
 pdfcreator={LaTeX via pandoc}}

\author{SCX}
\date{2026}

\begin{document}

\section{Proposition 3: State-Conditioned Weighting
Dominance}\label{proposition-3-state-conditioned-weighting-dominance}

\begin{quote}
2026-06-27: Reconstructed proof replacing the original weak version.
\end{quote}

\subsection{1 Overview}\label{overview}

The original Proposition 3 demonstrated via counterexample that global
fixed weights can fail. The present version replaces this with a
\textbf{rigorous inequality}: for any loss function satisfying mild
regularity, state-conditioned weighting achieves an expected risk no
larger than any global weighting scheme. The gap between the two
approaches is quantified and shown to be strictly positive precisely
when expert risk functions cross across states --- the exact condition
under which SCX provides benefit.

\begin{center}\rule{0.5\linewidth}{0.5pt}\end{center}

\subsection{2 Notation and Setup}\label{notation-and-setup}

\begin{longtable}[]{@{}
 >{\raggedright\arraybackslash}p{(\linewidth - 2\tabcolsep) * \real{0.4706}}
 >{\raggedright\arraybackslash}p{(\linewidth - 2\tabcolsep) * \real{0.5294}}@{}}
\toprule\noalign{}
\begin{minipage}[b]{\linewidth}\raggedright
Symbol
\end{minipage} & \begin{minipage}[b]{\linewidth}\raggedright
Meaning
\end{minipage} \\
\midrule\noalign{}
\endhead
\bottomrule\noalign{}
\endlastfoot
\(\mathcal{X}\) & Input space \\
\(\mathcal{Y}\) & Output space (could be labels or targets) \\
\(P_{X,Y}\) & Joint distribution over
\(\mathcal{X} \times \mathcal{Y}\) \\
\(\mathcal{S} = \{s_1,\dots,s_K\}\) & State partition of \(\mathcal{X}\)
induced by \(\Pi: \mathcal{X} \to \mathcal{S}\) \\
\(P(s) = P_X(\Pi^{-1}(s))\) & Probability mass of state \(s\) \\
\(\{f_1,\dots,f_M\}\) & Expert functions
\(f_m: \mathcal{X} \to \mathcal{Y}\) \\
\(\Delta^{M-1} = \{w \in \mathbb{R}^M_+ : \sum_{m=1}^M w_m = 1\}\) &
Probability simplex of ensemble weights \\
\(\ell: \mathcal{Y} \times \mathcal{Y} \to \mathbb{R}^+\) & Loss
function \\
\(f_w(x) = \sum_{m=1}^M w_m f_m(x)\) & Weighted ensemble prediction \\
\end{longtable}

\textbf{Assumption 1 (Lipschitz continuity).} The loss function
\(\ell(\cdot, y)\) is \(L\)-Lipschitz continuous in its first argument
for every \(y \in \mathcal{Y}\):

\[|\ell(y_1, y) - \ell(y_2, y)| \leq L \cdot \|y_1 - y_2\|, \quad \forall y_1, y_2 \in \mathcal{Y}, \; y \in \mathcal{Y}\]

This is satisfied by common losses: absolute error (\(L=1\)), squared
error on bounded domains (\(L = 2\sup|y|\)), hinge loss (\(L=1\)), and
cross-entropy with bounded logits.

\textbf{Assumption 2 (Risk finiteness).} For every expert \(m\) and
every state \(s\), the conditional risk
\(R_m(s) = \mathbb{E}_{X,Y \sim P_{X,Y|s}}[\ell(f_m(X), Y)]\) is finite.

\textbf{Assumption 3 (State coverage).} \(P(s) > 0\) for each
\(s \in \mathcal{S}\).

\begin{center}\rule{0.5\linewidth}{0.5pt}\end{center}

\subsection{3 Problem Formulation}\label{problem-formulation}

Define the \textbf{state-conditioned optimal risk}:

\[R_{\text{SC}} = \mathbb{E}_{s \sim P_S}\Bigl[\min_{w \in \Delta^{M-1}} \mathbb{E}_{x,y \sim P_{X,Y|s}}\bigl[\ell(f_w(x), y)\bigr]\Bigr]\]

This is the risk achieved by selecting, for each state \(s\), the
ensemble weights \(w(s)\) that minimize the conditional risk within that
state, then averaging over states.

Define the \textbf{globally optimal risk}:

\[R_{\text{Global}} = \min_{w \in \Delta^{M-1}} \mathbb{E}_{x,y \sim P_{X,Y}}\bigl[\ell(f_w(x), y)\bigr]\]

This is the risk achieved by the single best global weighting scheme.

Our goal is to prove \(R_{\text{SC}} \leq R_{\text{Global}}\) and to
characterize when the inequality is strict.

\begin{center}\rule{0.5\linewidth}{0.5pt}\end{center}

\subsection{4 Theorem 2: State-Conditioned
Dominance}\label{theorem-2-state-conditioned-dominance}

\subsubsection{4.1 Statement}\label{statement}

Under Assumptions 1-3:

\[R_{\text{SC}} \leq R_{\text{Global}}\]

Equivalently:

\[\mathbb{E}_{s}\Bigl[\min_{w} \mathbb{E}_{x,y|s}\bigl[\ell(f_w(x), y)\bigr]\Bigr] \leq \min_{w} \mathbb{E}_{x,y}\bigl[\ell(f_w(x), y)\bigr]\]

\subsubsection{4.2 Proof}\label{proof}

\textbf{Step 1: Pointwise inequality.} For any fixed global weight
\(w_0 \in \Delta^{M-1}\) and any state \(s \in \mathcal{S}\), consider
the conditional risk under \(w_0\):

\[\mathcal{R}(s, w_0) = \mathbb{E}_{x,y \sim P_{X,Y|s}}\bigl[\ell(f_{w_0}(x), y)\bigr]\]

The state-conditioned approach can, at minimum, match \(w_0\) in state
\(s\) by also choosing \(w_0\). Thus:

\[\min_{w \in \Delta^{M-1}} \mathcal{R}(s, w) \leq \mathcal{R}(s, w_0)\]

This holds for every \(s \in \mathcal{S}\) and every
\(w_0 \in \Delta^{M-1}\).

\textbf{Step 2: Expectation preserves the inequality.} Taking
expectation over \(s \sim P_S\) on both sides:

\[\mathbb{E}_s\Bigl[\min_{w} \mathcal{R}(s, w)\Bigr] \leq \mathbb{E}_s\bigl[\mathcal{R}(s, w_0)\bigr]\]

The right-hand side equals the global risk under weight \(w_0\):

\[\mathbb{E}_s\bigl[\mathcal{R}(s, w_0)\bigr] = \sum_{s \in \mathcal{S}} P(s) \cdot \mathbb{E}_{x,y|s}\bigl[\ell(f_{w_0}(x), y)\bigr] = \mathbb{E}_{x,y}\bigl[\ell(f_{w_0}(x), y)\bigr]\]

\textbf{Step 3: Minimizing over \(w_0\).} Since the inequality holds for
every \(w_0\), we can take the minimum over \(w_0\) on the right-hand
side, obtaining:

\[\mathbb{E}_s\Bigl[\min_{w} \mathcal{R}(s, w)\Bigr] \leq \min_{w_0} \mathbb{E}_{x,y}\bigl[\ell(f_{w_0}(x), y)\bigr]\]

which is precisely \(R_{\text{SC}} \leq R_{\text{Global}}\). \(\square\)

\subsubsection{4.3 Alternative Proof via Jensen's
Inequality}\label{alternative-proof-via-jensens-inequality}

The same result follows from the concavity of the minimum-over-weights
functional. For each \(s\), define:

\[\phi(s) = \min_{w \in \Delta^{M-1}} \mathbb{E}_{x,y|s}\bigl[\ell(f_w(x), y)\bigr]\]

The function \(\phi\) is a pointwise minimum of functions affine in the
conditional distribution \(P_{X,Y|s}\), making \(\phi\) a
\textbf{concave functional} of \(P_{X,Y|s}\). More concretely, for any
fixed \(w\), the map \(\psi_w(P) = \mathbb{E}_{P}[\ell(f_w(X), Y)]\) is
linear in \(P\), and \(\phi(P) = \min_w \psi_w(P)\) is the lower
envelope of linear functions, hence concave.

By Jensen's inequality for concave functionals:

\[\phi\bigl(\mathbb{E}_s[P_{X,Y|s}]\bigr) \geq \mathbb{E}_s\bigl[\phi(P_{X,Y|s})\bigr]\]

Since \(\mathbb{E}_s[P_{X,Y|s}] = P_{X,Y}\) (the marginal distribution),
we have:

\[\min_{w} \mathbb{E}_{x,y}\bigl[\ell(f_w(x), y)\bigr] = \phi(P_{X,Y}) \geq \mathbb{E}_s\bigl[\phi(P_{X,Y|s})\bigr] = \mathbb{E}_s\Bigl[\min_{w} \mathbb{E}_{x,y|s}\bigl[\ell(f_w(x), y)\bigr]\Bigr]\]

which is \(R_{\text{Global}} \geq R_{\text{SC}}\). The Jensen
perspective reveals that the inequality is structural: it is a direct
consequence of the concavity of the min-operator in the distribution
argument. \(\square\)

\begin{center}\rule{0.5\linewidth}{0.5pt}\end{center}

\subsection{5 The Gap and When It
Vanishes}\label{the-gap-and-when-it-vanishes}

\subsubsection{5.1 Definition of the Gap}\label{definition-of-the-gap}

Define the \textbf{state-conditioning advantage}:

\[\Delta = R_{\text{Global}} - R_{\text{SC}} \geq 0\]

Theorem 2 guarantees \(\Delta \geq 0\). We now characterize when
\(\Delta = 0\) and when \(\Delta > 0\).

\subsubsection{5.2 Theorem 3: Zero Gap
Condition}\label{theorem-3-zero-gap-condition}

\[\Delta = 0 \quad \Longleftrightarrow \quad \exists w^* \in \Delta^{M-1} \text{ such that } \forall s \in \mathcal{S}: w^* \in \arg\min_{w} \mathcal{R}(s, w)\]

That is, the gap is zero if and only if there exists a single global
weight vector that is simultaneously optimal in every state. When this
holds, the state-conditioned approach offers no benefit over the best
global weighting.

\subsubsection{5.3 Proof}\label{proof-1}

\textbf{(\(\Leftarrow\)).} If \(w^*\) is optimal in every state, then
\(\min_w \mathcal{R}(s, w) = \mathcal{R}(s, w^*)\) for each \(s\).
Taking expectations:

\[R_{\text{SC}} = \mathbb{E}_s[\mathcal{R}(s, w^*)] = R_{\text{Global}}\]

so \(\Delta = 0\).

\textbf{(\(\Rightarrow\)).} Suppose \(\Delta = 0\), i.e.,
\(R_{\text{SC}} = R_{\text{Global}}\). Let \(w^*\) be any global
minimizer, so \(R_{\text{Global}} = \mathbb{E}_s[\mathcal{R}(s, w^*)]\).
Since \(\min_w \mathcal{R}(s, w) \leq \mathcal{R}(s, w^*)\) for all
\(s\), we have:

\[0 = R_{\text{Global}} - R_{\text{SC}} = \mathbb{E}_s\bigl[\mathcal{R}(s, w^*) - \min_w \mathcal{R}(s, w)\bigr]\]

The integrand is everywhere non-negative (by the optimality of
\(\min_w\)), and the expectation is zero. Therefore the integrand must
be zero almost surely under \(P_S\):

\[\mathcal{R}(s, w^*) - \min_w \mathcal{R}(s, w) = 0 \quad \text{for } P_S\text{-almost every } s\]

Since \(P(s) > 0\) for all \(s\) (Assumption 3), this holds for every
\(s \in \mathcal{S}\), establishing that \(w^*\) is optimal in every
state. \(\square\)

\begin{center}\rule{0.5\linewidth}{0.5pt}\end{center}

\subsection{6 The Role of Risk
Crossing}\label{the-role-of-risk-crossing}

\subsubsection{6.1 Definition of Risk
Crossing}\label{definition-of-risk-crossing}

\textbf{Definition 4 (Risk crossing).} Expert risk functions exhibit
\textbf{crossing} across states if there exist two states
\(s_1, s_2 \in \mathcal{S}\) and two experts \(a, b \in \mathcal{M}\)
such that:

\[R_a(s_1) < R_b(s_1) \quad \text{and} \quad R_a(s_2) > R_b(s_2)\]

where \(R_m(s) = \mathbb{E}_{x,y|s}[\ell(f_m(x), y)]\) is the risk of
expert \(m\) alone (not the ensemble).

\subsubsection{6.2 Theorem 4: Crossing Implies Positive
Gap}\label{theorem-4-crossing-implies-positive-gap}

If the expert risk functions cross, and the loss \(\ell\) is strictly
convex in its first argument, then \(\Delta > 0\).

\subsubsection{6.3 Proof}\label{proof-2}

\textbf{Step 1: Crossing forces different optimal weights.} Under risk
crossing, expert \(a\) outperforms expert \(b\) in state \(s_1\), and
\(b\) outperforms \(a\) in \(s_2\). For strictly convex \(\ell\), the
optimal ensemble weight vector in each state will favor the better
expert:

\[\arg\min_w \mathcal{R}(s_1, w) \neq \arg\min_w \mathcal{R}(s_2, w)\]

To see this formally, consider the optimal simplex weights for two
experts (\(M=2\)). The ensemble prediction is
\(f_w(x) = w_1 f_1(x) + w_2 f_2(x)\) with \(w_1 + w_2 = 1\). With
strictly convex \(\ell\), the function
\(w_1 \mapsto \mathcal{R}(s, w_1)\) is strictly convex in \(w_1\). The
first-order condition gives the unique optimal weight \(w_1^*(s)\). In
state \(s_1\), since \(R_1(s_1) < R_2(s_1)\), the optimal weight
satisfies \(w_1^*(s_1) > 0.5\) (expert 1 is favored). In state \(s_2\),
since \(R_1(s_2) > R_2(s_2)\), the optimal weight satisfies
\(w_1^*(s_2) < 0.5\) (expert 2 is favored). Hence
\(w^*(s_1) \neq w^*(s_2)\).

For \(M > 2\), a similar argument holds: the crossing condition implies
that the set \(\arg\min_w \mathcal{R}(s, w)\) differs between \(s_1\)
and \(s_2\) because the ranking of risk values \(R_m(s)\) changes,
altering the KKT conditions of the convex optimization.

\textbf{Step 2: No single weight is optimal everywhere.} Since
\(w^*(s_1) \neq w^*(s_2)\), there cannot exist a single \(w^*\) that is
simultaneously optimal in both states. By Theorem 3, this implies
\(\Delta > 0\).

\textbf{Step 3: Strictness argument.} Since the risk functions cross,
\(\Delta\) cannot be zero. By Theorem 2, \(\Delta \geq 0\), so
\(\Delta > 0\) necessarily. \(\square\)

\subsubsection{6.4 Remarks}\label{remarks}

\begin{itemize}
\item
 \textbf{Non-strictly convex losses.} For losses that are convex but
 not strictly convex (e.g., absolute error with potentially flat
 minima), crossing implies \(\Delta \geq 0\), but the inequality may
 not be strict. In this case, the optimal weight set may contain
 multiple values, and it is possible (though measure-zero) that the
 same weight is optimal in both crossing states despite the crossing.
 Generic position (non-degenerate risks) restores strict positivity.
\item
 \textbf{The converse.} \(\Delta > 0\) does not necessarily imply
 pairwise risk crossing as defined in Definition 4. The gap can be
 positive even without pairwise crossing if three or more experts
 interact such that no single weight is optimal everywhere, but any two
 experts individually appear co-monotonic. The gap condition of Theorem
 3 --- existence of a universally optimal weight vector --- is the
 precise characterization.
\end{itemize}

\begin{center}\rule{0.5\linewidth}{0.5pt}\end{center}

\subsection{7 Practical Corollaries}\label{practical-corollaries}

\subsubsection{7.1 Corollary 1: Gap Lower
Bound}\label{corollary-1-gap-lower-bound}

Under the Lipschitz assumption (Assumption 1), the advantage \(\Delta\)
is bounded below by the variation in optimal weights:

\[\Delta \geq \frac{L}{2} \cdot \mathbb{E}_s\bigl[\|w^*(s) - w^*_{\text{global}}\|_1\bigr] \cdot \bigl(\min_{s} \min_{m} \mathbb{E}_{x|s}[\|f_m(x)\|]\bigr)\]

where \(w^*(s)\) is the optimal weight in state \(s\) and
\(w^*_{\text{global}}\) is the global optimal weight.

\textbf{Proof sketch.} For Lipschitz \(\ell\), the risk difference
\(|\mathcal{R}(s, w) - \mathcal{R}(s, w')|\) can be bounded by
\(L \cdot \|w - w'\|_1 \cdot \mathbb{E}_{x|s}[\max_m \|f_m(x)\|]\).
Applying this to the gap definition yields the bound. \(\square\)

\subsubsection{7.2 Corollary 2: Finite-Sample
Bound}\label{corollary-2-finite-sample-bound}

Let \(\hat{R}_m(s)\) be empirical risk estimates based on \(n_s\)
samples in state \(s\). The empirical advantage:

\[\hat{\Delta} = \min_{w} \sum_{s} P(s) \hat{\mathcal{R}}(s, w) - \sum_{s} P(s) \min_{w} \hat{\mathcal{R}}(s, w)\]

satisfies \(\hat{\Delta} \to \Delta\) almost surely as
\(\min_s n_s \to \infty\), and with probability \(1-\eta\):

\[|\hat{\Delta} - \Delta| \leq 2L_{\max} \sum_{s} P(s) \sqrt{\frac{2\log(4MK/\eta)}{n_s}}\]

This enables practitioners to compute statistically valid estimates of
the benefit of state conditioning.

\subsubsection{7.3 Corollary 3: Monotonicity in State
Granularity}\label{corollary-3-monotonicity-in-state-granularity}

Let \(\mathcal{S}_1\) and \(\mathcal{S}_2\) be two state partitions such
that \(\mathcal{S}_2\) is a refinement of \(\mathcal{S}_1\) (every state
in \(\mathcal{S}_1\) is a union of states in \(\mathcal{S}_2\)). Then:

\[\Delta_{\mathcal{S}_2} \geq \Delta_{\mathcal{S}_1}\]

The advantage of state conditioning never decreases under finer state
partitioning.

\textbf{Proof.} State-conditioned risk with partition \(\mathcal{S}_2\)
is:

\[R_{\text{SC}}(\mathcal{S}_2) = \mathbb{E}_{s_2}\bigl[\min_w \mathbb{E}_{x,y|s_2}[\ell]\bigr]\]

By the same Jensen/min inequality,
\(R_{\text{SC}}(\mathcal{S}_2) \leq R_{\text{SC}}(\mathcal{S}_1)\)
(finer partition gives lower or equal risk). The global risk
\(R_{\text{Global}}\) is the same for both partitions, so
\(\Delta_{\mathcal{S}_2} = R_{\text{Global}} - R_{\text{SC}}(\mathcal{S}_2) \geq R_{\text{Global}} - R_{\text{SC}}(\mathcal{S}_1) = \Delta_{\mathcal{S}_1}\).
\(\square\)

This monotonicity formalizes the intuition that more states (finer
granularity) can only improve the benefit of state conditioning, with
the limit \(|\mathcal{S}| = |\mathcal{X}|\) achieving the
information-theoretic upper bound.

\begin{center}\rule{0.5\linewidth}{0.5pt}\end{center}

\subsection{8 Connection to Jensen's Inequality: A Deeper
View}\label{connection-to-jensens-inequality-a-deeper-view}

The inequality \(R_{\text{SC}} \leq R_{\text{Global}}\) admits a
revealing interpretation through Jensen's inequality and the concavity
of the minimum operator.

Define a function \(\varphi\) on the space of probability distributions
over \(\mathcal{X} \times \mathcal{Y}\):

\[\varphi(P) = \min_{w \in \Delta^{M-1}} \mathbb{E}_{(x,y) \sim P}\bigl[\ell(f_w(x), y)\bigr]\]

For any two distributions \(P_1, P_2\) and \(\lambda \in [0,1]\):

\[\varphi(\lambda P_1 + (1-\lambda) P_2) = \min_w \bigl[\lambda \mathbb{E}_{P_1}[\ell(f_w)] + (1-\lambda) \mathbb{E}_{P_2}[\ell(f_w)]\bigr]\]
\[\geq \lambda \min_w \mathbb{E}_{P_1}[\ell(f_w)] + (1-\lambda) \min_w \mathbb{E}_{P_2}[\ell(f_w)]\]
\[= \lambda \varphi(P_1) + (1-\lambda) \varphi(P_2)\]

where the inequality holds because the minimum of a sum is at least the
sum of the minima (a classic inequality:
\(\min_x (a(x)+b(x)) \geq \min_x a(x) + \min_x b(x)\)). This establishes
that \(\varphi\) is \textbf{concave}.

Jensen's inequality for concave \(\varphi\) states:

\[\varphi(\mathbb{E}[P]) \geq \mathbb{E}[\varphi(P)]\]

Applying this with the random distribution \(P_{X,Y|S}\) (the
conditional distribution given state \(S\)):

\[\varphi(\mathbb{E}_S[P_{X,Y|S}]) \geq \mathbb{E}_S[\varphi(P_{X,Y|S})]\]

Since \(\mathbb{E}_S[P_{X,Y|S}] = P_{X,Y}\), this yields:

\[\min_w \mathbb{E}_{X,Y}[\ell(f_w)] \geq \mathbb{E}_S\Bigl[\min_w \mathbb{E}_{X,Y|S}[\ell(f_w)]\Bigr]\]

i.e., \(R_{\text{Global}} \geq R_{\text{SC}}\), exactly the result of
Theorem 2.

The concavity perspective reveals that state conditioning exploits the
\textbf{averaging-is-beneficial} property of concave functions: the
average of per-state minima is no worse than the minimum of the average
distribution.

\subsubsection{8.1 Entropy Interpretation of the
Gap}\label{entropy-interpretation-of-the-gap}

The state-conditioning advantage admits an information-theoretic
interpretation via entropy. Let \(w_{\text{global}} \in \Delta^{M-1}\)
denote the globally optimal weight vector, and let
\(w^*(s) \in \Delta^{M-1}\) denote the optimal weight vector for state
\(s\). Define the entropy of a weight distribution as
\(H(w) = -\sum_{m=1}^M w_m \log w_m\).

By concavity of entropy, for any mixture of distributions:
\[H\!\left(\sum_s P(s) \cdot w^*(s)\right) \geq \sum_s P(s) \cdot H(w^*(s)) = \mathbb{E}_s[H(w^*(s))].\]

Under the typical condition that \(w_{\text{global}}\) approximates the
state-average of per-state optima (i.e.,
\(w_{\text{global}} \approx \mathbb{E}_s[w^*(s)]\), which holds exactly
when the global risk minimizer lies in the convex hull of per-state
minimizers), we obtain:

\[\boxed{\;H(w_{\text{global}}) \geq \mathbb{E}_s[H(w^*(s))]\;}.\]

The inequality is oriented \(\geq\), not \(\leq\): the global weight
vector has entropy \textbf{at least} the expected per-state entropy.
This reflects that state-conditioned weights are typically more
concentrated (lower entropy per state) because each state's optimal
weights favor the experts that perform best in that state, while the
global weights must compromise across states, producing a more dispersed
(higher entropy) distribution. The gap
\(H(w_{\text{global}}) - \mathbb{E}_s[H(w^*(s))] \geq 0\) quantifies how
much ``sharper'' the per-state weighting is relative to the global
compromise --- it is an information-theoretic measure of the benefit of
state conditioning.

\subsubsection{8.2 Correlation Condition for Weighting
Improvement}\label{correlation-condition-for-weighting-improvement}

In analyzing whether a weight update \(w \to w'\) reduces the ensemble
risk, one encounters the cross-term
\(\mathbb{E}[(\ell - \hat{R})(w - w')]\), where \(\ell\) is the
per-sample loss, \(\hat{R}\) is the estimated risk, and \(w, w'\) are
weight vectors. The sign of this expectation determines whether the
weight change is aligned with error reduction.

The original claim \(\mathbb{E}[(\ell - \hat{R})(w - w')] \geq 0\) does
\textbf{not} hold in general. It is equivalent to the condition that
weight adjustments are positively correlated with error residuals ---
increasing weight on experts whose current loss exceeds their estimated
risk, and decreasing weight otherwise. However, without further
structure, the sign of this covariance is unconstrained:

\[\mathbb{E}[(\ell - \hat{R})(w - w')] = \operatorname{Cov}(\ell - \hat{R},\; w - w') + \mathbb{E}[\ell - \hat{R}] \cdot \mathbb{E}[w - w'].\]

\textbf{Revised statement (working conjecture).} This inequality holds
under the additional assumption
\(\operatorname{Cov}(\ell - \hat{R},\; w - w') = 0\). Without this, the
sign is not guaranteed. We adopt this as a \textbf{working conjecture}:
for well-specified SCX state partitions and experts with conditionally
independent errors (Assumption A2), the error residual
\(\ell - \hat{R}\) is approximately uncorrelated with the weight
adjustment direction \(w - w'\), making the non-negativity condition
plausible. A rigorous proof of this uncorrelatedness remains open and
likely requires stronger conditions on the state partition and expert
error structure.

\begin{center}\rule{0.5\linewidth}{0.5pt}\end{center}

\subsection{9 Implications for SCX}\label{implications-for-scx}

\begin{enumerate}
\def\labelenumi{\arabic{enumi}.}
\item
 \textbf{Theoretical guarantee.} Theorem 2 establishes a rigorous
 dominance result: state-conditioned weighting is never worse than
 global weighting. The proof requires only the existence of a
 well-defined state partition and does not depend on distributional
 assumptions beyond finiteness.
\item
 \textbf{Gap as a diagnostic.} The advantage \(\Delta\) quantifies the
 value of state conditioning. When \(\Delta\) is small, the state
 partition may be too coarse or the experts may not have complementary
 strengths. SCX's state discovery module should aim to maximize
 \(\Delta\).
\item
 \textbf{Risk crossing is the engine.} Theorem 4 confirms that the
 crossing of expert risk functions --- the fundamental phenomenon that
 SCX exploits --- is necessary and sufficient for a strictly positive
 advantage (under strict convexity).
\item
 \textbf{Granularity trade-off.} Corollary 3 shows that finer states
 always help in expectation, but at the cost of increased estimation
 variance (Corollary 2). SCX's adaptive state discovery balances this
 bias-variance trade-off.
\item
 \textbf{Unified framework.} The Jensen perspective (Section 8) unifies
 the SCX weighting formula with information-theoretic principles: the
 gap \(\Delta\) equals the expected reduction in risk from conditioning
 on the state variable, analogous to the reduction in variance in
 Rao-Blackwellization.
\end{enumerate}

\begin{center}\rule{0.5\linewidth}{0.5pt}\end{center}

\subsection{References}\label{references}

\begin{enumerate}
\def\labelenumi{\arabic{enumi}.}
\tightlist
\item
 SCX theoretical framework. \texttt{01\_SCX\_Core Framework\_Mathematical Analysis.md}
 Section 2.3
\item
 SCX Mathematical Foundations. \texttt{05\_Mathematical Roots and Proofs.md}
 Proposition 3
\item
 Boyd, S. \& Vandenberghe, L. (2004). \emph{Convex Optimization}.
 Cambridge University Press. {[}Chapter 3: concavity of pointwise
 minimum{]}
\item
 Rao, C. R. (1945). Information and accuracy attainable in the
 estimation of statistical parameters. \emph{Bulletin of the Calcutta
 Mathematical Society}, 37, 81-91. {[}Rao-Blackwell theorem{]}
\item
 Tsybakov, A. B. (2004). Optimal aggregation of classifiers in
 statistical learning. \emph{Annals of Statistics}, 32(1), 135-166.
\end{enumerate}

\begin{center}\rule{0.5\linewidth}{0.5pt}\end{center}

\subsection{Changelog}\label{changelog}

\begin{longtable}[]{@{}
 >{\raggedright\arraybackslash}p{(\linewidth - 6\tabcolsep) * \real{0.1875}}
 >{\raggedright\arraybackslash}p{(\linewidth - 6\tabcolsep) * \real{0.2500}}
 >{\raggedright\arraybackslash}p{(\linewidth - 6\tabcolsep) * \real{0.2500}}
 >{\raggedright\arraybackslash}p{(\linewidth - 6\tabcolsep) * \real{0.3125}}@{}}
\toprule\noalign{}
\begin{minipage}[b]{\linewidth}\raggedright
Date
\end{minipage} & \begin{minipage}[b]{\linewidth}\raggedright
Defect
\end{minipage} & \begin{minipage}[b]{\linewidth}\raggedright
Change
\end{minipage} & \begin{minipage}[b]{\linewidth}\raggedright
Severity
\end{minipage} \\
\midrule\noalign{}
\endhead
\bottomrule\noalign{}
\endlastfoot
2026-06-28 & B4 & \textbf{Added entropy inequality} (§8.1):
\(H(w_{\text{global}}) \geq \mathbb{E}_s[H(w^*(s))]\) (correct
orientation: global entropy is \textbf{at least} expected per-state
entropy). The original claim
\(H(w_{\text{global}}) \leq \mathbb{E}_s[H(w^*(s))]\) was reversed ---
by concavity of entropy, the inequality runs \(\geq\), not \(\leq\). &
MAJOR \\
2026-06-28 & B5 & \textbf{Corrected correlation claim} (§8.2): Replaced
the unqualified \(\mathbb{E}[(\ell - \hat{R})(w - w')] \geq 0\) with a
working conjecture that requires
\(\operatorname{Cov}(\ell - \hat{R}, w - w') = 0\). Without this
additional assumption, the sign is not guaranteed. & MAJOR \\
\end{longtable}

\end{document}
